\section{Norm-visible central matrix cones}\label{norm:section}

The construction in this section detects a prescribed positive element
uniformly in the irreducible representations at each coordinate.  Its
building blocks are finite word spaces at pure states.  Lemma~\ref{lem:shrinking}
lets us separate their hereditary supports while retaining all the
representations that must be detected.

\begin{proposition}\label{prop:norm}
Suppose that \(A\) is separable and satisfies \textup{(C)}.  For every positive contraction \(x\in A\), every
\(\gamma>0\), and every integer \(k\ge2\), there are \cpc maps
\(\nu_n:M_k\to A\) such that
\[
 \nu:M_k\longrightarrow \Aomega\cap A',
 \qquad \nu(z)=[(\nu_n(z))_n],
\]
is order zero, and there is a number \(\kappa>0\) for which
\begin{equation}\label{norm:coordinate-visibility}
 \begin{gathered}
 \text{for every }n\in\N\text{ and every irreducible representation }\sigma
 \text{ of }A,\\
 \|\sigma(x)\|\ge\gamma
 \ \Longrightarrow\
 \|\sigma(x\nu_n(1))\|\ge\kappa .
 \end{gathered}
\end{equation}
One may take \(\kappa=\gamma^2/32\).  In particular, the estimate holds on
one \(\omega\)-large set of coordinates simultaneously for every
irreducible representation, including representations chosen separately
at each coordinate.
\end{proposition}

We first establish the uniformity and compression facts used in the proof.
All support projections below are taken in the bidual, or in the universal
representation, and each representation is extended normally to the
bidual when support estimates are used.

\subsection{Finite covers and pure-state compression}

For \(x\in A_+\) and \(\gamma>0\), write
\[
 \mathcal K(x,\gamma)
   =\{\sigma:\sigma\text{ is an irreducible representation of }A,\
                    \|\sigma(x)\|\ge\gamma\}.
\]
There is no compactness assumption on the primitive ideal space.

\begin{lemma}\label{norm:finite-cover}
Let \(x\in A_+\), \(\gamma>0\), and let
\((s_\alpha)_{\alpha\in\Lambda}\) be positive elements of \(A\).
Suppose that \(\mathcal K(x,\gamma)\ne\varnothing\) and that, for
every \(\sigma\in\mathcal K(x,\gamma)\), some
\(\sigma(s_\alpha)\) is nonzero.  Then there are
\(\alpha_1,\ldots,\alpha_m\) and \(t>0\) such that
\[
 \sigma\in\mathcal K(x,\gamma)
 \quad\Longrightarrow\quad
 \max_{1\le i\le m}\|\sigma(s_{\alpha_i})\|>t.
\]
Thus the nonvanishing loci of the cuts
\((s_{\alpha_i}-t)_+\) still cover \(\mathcal K(x,\gamma)\).
For a finite covering family of positive contractions, continuous
functional calculus can replace its members by positive contractions
with the same supports and with maximum represented norm equal to one
at every \(\sigma\in\mathcal K(x,\gamma)\).
\end{lemma}

\begin{proof}
Let \(I\) be the closed ideal generated by all the \(s_\alpha\), and let
\(q:A\to A/I\) be the quotient map.  We claim that
\(\|q(x)\|<\gamma\).  Otherwise, a pure state of the unitization of
\(A/I\) can be chosen to take the value \(\|q(x)\|\) on \(q(x)\).
For completeness, the states taking that value form a nonempty compact
face of the state space: nonemptiness follows by extending the
norm-attaining state of the commutative algebra generated by \(q(x)\)
and the unit.  An extreme point of that face is a pure state.
Its value on \(q(x)\) is positive, so its irreducible representation
restricts nontrivially to \(A/I\).  Pulling that representation back to
\(A\) gives an element of \(\mathcal K(x,\gamma)\) annihilating every
\(s_\alpha\), a contradiction.

The strict quotient-norm inequality permits a finite sum
\[
 y=\sum_{j=1}^r u_j s_{\alpha(j)}v_j,\qquad
 u_j,v_j\in\widetilde A,\qquad \|x-y\|<\gamma .
\]
Indeed, these finite sums are dense in \(I\), and the distance from
\(x\) to \(I\) is strictly less than \(\gamma\).  Choose \(t>0\) so
small that
\[
 y_t=\sum_{j=1}^r u_j(s_{\alpha(j)}-t)_+v_j
 \quad\text{satisfies}\quad \|x-y_t\|<\gamma .
\]
If a representation annihilates every cut appearing in this sum, then
\(\|\sigma(x)\|\le\|x-y_t\|<\gamma\).  This proves the first assertions.

For the last assertion apply what was just proved to the given finite
family and put \(f_t(u)=\min\{u/t,1\}\) for \(u\ge0\).
The function \(f_t\) vanishes only at zero, so it preserves supports.
Whenever \(\|\sigma(s_i)\|>t\), functional calculus gives
\(\|\sigma(f_t(s_i))\|=1\).
\end{proof}

The following compression is the pure-state excision principle of
Akemann, Anderson, and Pedersen \cite[Proposition~2.2]{AAP1986}.
We include a proof using Kadison's transitivity theorem and the
pure-state kernel identity \cite[Theorem~1 and Corollary~2]{Kadison1957}.

\begin{lemma}[Pure-state compression]\label{norm:excision}
Let \(\lambda\) be a pure state of a unital \(C^*\)-algebra \(D\), with
irreducible GNS representation \((\pi,\mathcal H,\xi)\).
Given a finite subset \(\mathcal E\subset D\) and \(\epsilon>0\), there
is a positive contraction \(q\in D\) such that
\[
 \pi(q)\xi=\xi,\qquad
 \|qyq-\lambda(y)q^2\|<\epsilon\quad(y\in\mathcal E).
\]
\end{lemma}

\begin{proof}
Set
\[
 J=\{z\in D:\pi(z)\xi=0\},\qquad H=J\cap J^*.
\]
The space \(J\) is a closed left ideal, and \(H\) is a hereditary
\(C^*\)-subalgebra.  For the latter assertion, it is a closed
selfadjoint subalgebra, and if \(0\le a\le b\) with \(b\in H_+\), then
\(\lambda(a)\le\lambda(b)=0\), so \(\pi(a)\xi=0\).

We show that
\begin{equation}\label{norm:state-kernel}
 \ker\lambda=J+J^*.
\end{equation}
The inclusion from right to left is immediate.  For the converse,
let \(\lambda(y)=0\) and put \(\zeta=\pi(y)\xi\), which is orthogonal
to \(\xi\).  If \(\zeta=0\), then \(y\in J\).  Otherwise, on
\(K=\operatorname{span}\{\xi,\zeta\}\), define \(R\xi=\zeta\) and
\(R\zeta=0\).  The selfadjoint form of Kadison's Transitivity Theorem
\cite{Kadison1957} supplies selfadjoint \(h,g\in D\) whose restrictions
to \(K\) are respectively \(\operatorname{Re}R\) and
\(\operatorname{Im}R\).  Thus \(v=h+ig\) satisfies
\[
 \pi(v)\xi=\zeta,\qquad \pi(v^*)\xi=0.
\]
Consequently \(y-v\in J\) and \(v\in J^*\), proving
\eqref{norm:state-kernel}.

Let \((e_\alpha)\) be a positive contractive approximate unit for \(H\)
and put \(q_\alpha=1-e_\alpha\).  Each \(q_\alpha\) fixes \(\xi\).
For \(z\in J\), the element \(z^*z\) belongs to \(H\), and hence
\[
 \|zq_\alpha\|^2
   =\|q_\alpha z^*zq_\alpha\|\longrightarrow0.
\]
Its adjoint gives \(\|q_\alpha z^*\|\to0\).  Applying
\eqref{norm:state-kernel} to \(y-\lambda(y)1\) now gives
\[
 \|q_\alpha(y-\lambda(y)1)q_\alpha\|\longrightarrow0
 \qquad(y\in D).
\]
Choose one index sufficiently late for all the finitely many tests.
\end{proof}

We will apply excision to make finitely many word representatives
almost orthonormal on one hereditary support.  Their Gram matrix will
control the norm of the row formed from these representatives.  We
also estimate how that row intertwines the action of each prescribed
element below the outermost word level; a taper will later remove
that boundary level.

To construct this matrix form of the compression, fix a finite set
\(F\subset A\) of selfadjoint contractions, \(d\ge1\), and
\(0<\eta<1\).  Suppose \(\sigma\in\mathcal K(x,\gamma)\), where \(x\)
is a positive contraction.  Choose a unit vector \(\xi\) with
\[
 \lambda(x):=\langle\sigma(x)\xi,\xi\rangle>\gamma/2.
\]
Every nonzero vector in an irreducible representation is cyclic, so
this is a pure state of \(\widetilde A\).  Let \(\mathcal H_l\) be the
span of the vectors \(\sigma(w)\xi\), where \(w\) is a word of length
at most \(l\) in \(F\), with the empty word equal to \(1\).  These
finite-dimensional spaces form an increasing filtration
\[
 \C\xi=\mathcal H_0\subseteq\mathcal H_1
       \subseteq\cdots\subseteq\mathcal H_d.
\]
Choose an orthonormal basis \(\xi_1,\ldots,\xi_m\) adapted to this
filtration, with \(\xi_1=\xi\).  Write
\(\xi_j=\sigma(b_j)\xi\), where \(b_j\in\widetilde A\) is a finite
linear combination of the appropriate words; take \(b_1=1\).
Let \(l_j\in\{0,\ldots,d\}\) be the level of \(\xi_j\).  There is
exactly one level-zero vector.

For \(a\in F\), let \(M_a\in M_m\) be the matrix of its compression
to \(\mathcal H_d\) in this basis.  It is a selfadjoint contraction.
Moreover,
\begin{equation}\label{norm:banded-model}
 (M_a)_{ij}=0\quad\text{if }|l_i-l_j|>1.
\end{equation}
Indeed, \(a\mathcal H_l\subseteq\mathcal H_{l+1}\) for \(l<d\);
orthogonality to the filtration proves one half of the assertion,
and selfadjointness proves the other.  For \(l_j<d\), put
\[
 r_j(a)=ab_j-\sum_{l=1}^m b_l(M_a)_{lj}.
\]
Since \(a\xi_j\in\mathcal H_d\), these residuals satisfy
\(\sigma(r_j(a))\xi=0\).

\begin{lemma}\label{norm:prototype}
For this finite word model there is a positive contractive seed
\(s\in A\) with \(\sigma(s)\ne0\), whose support \(P=p_s\) satisfies
\begin{align}
 PxP&\ge(\gamma/4)P,                                      \label{norm:positive-compression}\\
 \big\|(Pb_i^*b_jP)_{i,j=1}^m-I_m\otimes P\big\|
      &\le\eta^2/4,                                     \label{norm:local-gram}\\
 \big\|(Pr_i(a)^*r_j(a)P)_{l_i,l_j<d}\big\|
      &\le\eta^2/4\qquad(a\in F).                        \label{norm:local-residual}
\end{align}
Every positive element in \(\her(s)\) has a support on which the
same three estimates hold.
\end{lemma}

\begin{proof}
Apply Lemma~\ref{norm:excision} in \(\widetilde A\) to \(x\), all
\(b_i^*b_j\), and all \(r_i(a)^*r_j(a)\).  Their state values are
respectively \(\lambda(x)>\gamma/2\), \(\delta_{ij}\), and zero.
Choose the excision error at most \(\gamma/16\) for \(x\), and at
most \(\eta^2/(16m)\) for each of the matrix entries.
Put
\[
 P_0=1_{(1/2,1]}(q),\qquad v=q^{-1}P_0
 \quad\text{in }\widetilde A^{**}.
\]
Here \(\|v\|\le2\), and
\[
 P_0(y-\lambda(y)1)P_0
       =v\big(qyq-\lambda(y)q^2\big)v.
\]
The estimate for \(x\) therefore gives
\(P_0xP_0\ge(\lambda(x)-\gamma/4)P_0\ge(\gamma/4)P_0\).
For the matrices, each entry has error at most
\(\eta^2/(4m)\).  An \(r\)-by-\(r\) operator matrix whose entries
have norm at most \(\delta\) has norm at most \(r\delta\), by the
Cauchy--Schwarz inequality for its row action.  Since each matrix
here has size at most \(m\), both matrix errors are at most
\(\eta^2/4\).

Choose a continuous function \(f:[0,1]\to[0,1]\) that is zero on
\([0,1/2]\) and has \(f(1)=1\), and set
\(s=f(q)xf(q)\).  This is a positive contraction in \(A\), because
\(A\) is an ideal in its unitization.  Also
\(\lambda(s)=\lambda(x)>0\), so \(\sigma(s)\ne0\), and \(p_s\le P_0\).
Compressing the preceding inequalities to \(p_s\) proves the
claims.  Further hereditary restriction only decreases the support
projection, so preserves all the estimates.
\end{proof}

\subsection{Separating the models}

The next observation records precisely why iterating
Lemma~\ref{lem:shrinking} does not lose the representation cover.

\begin{lemma}\label{norm:separated-cover}
Assume that $A$ satisfies \textup{(C)}. Suppose that positive
contractions \(s_1,\ldots,s_m\in A\) have
nonvanishing loci covering \(\mathcal K(x,\gamma)\).
For each pair \(i\ne j\), prescribe a finite set
\(\mathcal E_{ij}\subset\widetilde A\), and let \(\delta>0\).
There are positive contractions \(t_i\in\her(s_i)\), with supports
\(Q_i=p_{t_i}\), such that
\begin{align}
 \max_i\|\sigma(t_i)\|&=1
                 &&(\sigma\in\mathcal K(x,\gamma)),       \label{norm:retained-cover}\\
 \|Q_i z Q_j\|&\le\delta
                 &&(i\ne j,\ z\in\mathcal E_{ij}).        \label{norm:cross-compressions}
\end{align}
\end{lemma}

\begin{proof}
If \(\mathcal K(x,\gamma)\) is empty, take every \(t_i=0\).
Otherwise proceed as follows.
Let \(c<1\) be the constant in Lemma~\ref{lem:shrinking}.
Before each shrinking operation, apply the last assertion of
Lemma~\ref{norm:finite-cover} to the current finite family.
This leaves every support unchanged and ensures that each
representation in \(\mathcal K(x,\gamma)\) gives norm one to at
least one of the current seeds.

For one prescribed triple \(i,j,z\), let
\(\beta=\|p_{s_i}zp_{s_j}\|\), using the current seeds.  If
\(\beta=0\), nothing is needed.  Otherwise apply
Lemma~\ref{lem:shrinking} to the current seeds \(s_i,s_j\) and to
\(z/\beta\).  Their replacements have supports on which the
compression of \(z\) has norm at most \(c\beta\).
For any \(\sigma\in\mathcal K(x,\gamma)\), choose a current seed of
represented norm one.  If its index is outside \(\{i,j\}\), it
survives unchanged.  If its index is in \(\{i,j\}\), the
representation clause in Lemma~\ref{lem:shrinking} ensures that
at least one replacement is nonzero.  Thus the new family still
covers \(\mathcal K(x,\gamma)\), allowing the next
reparametrization.

All new supports are subprojections of their previous supports.
Consequently no previously obtained compression bound can increase.
For a fixed coefficient \(z\), finitely many repetitions suffice:
after \(r\) repetitions its bound is at most \(c^r\|z\|\).
There are only finitely many prescribed coefficients and pairs.
Treating them successively therefore gives
\eqref{norm:cross-compressions}.  A final support-preserving
reparametrization gives \eqref{norm:retained-cover}.
\end{proof}

\subsection{The row map and the taper}

\begin{lemma}\label{norm:finite-stage}
Assume that $A$ satisfies \textup{(C)}. Let \(x\in A\) be a positive
contraction, \(\gamma>0\), \(k\ge2\),
\(F\subset A\) a finite set of selfadjoint contractions,
\(d\ge1\), and \(0<\eta<1\).
There is a \cpc map \(V:M_k\to A\) such that
\begin{align}
 \|V(z)V(w)\|&\le\eta^2
   &&(z,w\in(M_k)_+,\ \|z\|,\|w\|\le1,\ zw=0),            \label{norm:order-zero-error}\\
 \|[a,V(z)]\|&\le2\eta+2/d
   &&(a\in F,\ \|z\|\le1),                               \label{norm:central-error}\\
 \|\sigma(xV(1))\|&\ge\gamma^2/32
   &&(\sigma\in\mathcal K(x,\gamma)).                    \label{norm:finite-visibility}
\end{align}
\end{lemma}

\begin{proof}
If \(\mathcal K(x,\gamma)\) is empty, take \(V=0\).  Otherwise
construct the preceding word model and the seed in
Lemma~\ref{norm:prototype} for each representation in that set.
The seeds cover the set.  Lemma~\ref{norm:finite-cover} therefore
retains finitely many models, indexed by \(i=1,\ldots,m\).
Write \(n_i\) for their basis sizes,
\(b_j^{(i)}\) for their basis representatives, \(l_j^{(i)}\) for
their levels, and \(M_a^{(i)}\) and \(r_j^{(i)}(a)\) for their
model matrices and residuals.  Set \(N=\sum_i n_i\).

Apply Lemma~\ref{norm:separated-cover} with
\(\delta=\eta^2/(2N)\) to all coefficients
\[
 (b_j^{(i)})^*b_l^{(h)}
 \quad\text{and}\quad
 (r_j^{(i)}(a))^*r_l^{(h)}(a)
 \qquad(i\ne h,\ a\in F),
\]
where the residual indices have levels less than \(d\).
Let \(s_i\) denote the resulting seeds and \(P_i=p_{s_i}\).
All the estimates of Lemma~\ref{norm:prototype} remain valid.
The number of models and their sizes have already been fixed
before this cross-compression tolerance is chosen.

By Lemma~\ref{lem:placement}, choose \cpc order-zero maps
\(\phi_i:M_k\to\her(s_i)\) such that
\begin{equation}\label{norm:cone-norm-one}
 \|\sigma(s_i)\|=1\quad\Longrightarrow\quad
 \|\sigma(\phi_i(1))\|=1.
\end{equation}
In fact it suffices to apply that Lemma to the fixed cut
\((s_i-1/2)_+\).  In the universal representation on
\(\mathcal H\), consider
\[
 \mathcal L=\bigoplus_{i=1}^m\bigoplus_{j=1}^{n_i}P_i\mathcal H
 \quad\text{and}\quad
 B:\mathcal L\to\mathcal H,\qquad
 B(\zeta_{ij})=\sum_{i,j}b_j^{(i)}\zeta_{ij}.
\]
The row \(B\) is only an operator on these support spaces.
Its Gram matrix has \((ij,hl)\)-entry
\(P_i(b_j^{(i)})^*b_l^{(h)}P_h\).
The diagonal blocks differ from the corresponding identity blocks
by at most \(\eta^2/4\).  The off-diagonal matrix has \(N\) rows
and columns with entries bounded by \(\delta\), so has norm at most
\(N\delta=\eta^2/2\).  In particular,
\begin{equation}\label{norm:row-gram}
 \|B^*B-1_{\mathcal L}\|\le\eta^2,\qquad
 \|B\|\le(1+\eta^2)^{1/2}.
\end{equation}

Let
\[
 M_a=\bigoplus_i(M_a^{(i)}\otimes1_{P_i\mathcal H})
       \quad\text{on }\mathcal L,
\]
and let \(Q\) be the projection onto the slots of level less than
\(d\).  The columns of \((aB-BM_a)Q\) are
\(r_j^{(i)}(a)P_i\).  Its squared norm is the norm of their Gram
matrix.  The local residual estimates and the same off-diagonal
estimate therefore give
\begin{equation}\label{norm:row-residual}
 \|(aB-BM_a)Q\|\le\eta\qquad(a\in F).
\end{equation}

For \(z\in M_k\), define operators on \(\mathcal L\) by
\[
 D_z=\bigoplus_{i,j}\phi_i(z)|_{P_i\mathcal H},\qquad
 T=\bigoplus_{i,j}(1-l_j^{(i)}/d)1_{P_i\mathcal H},\qquad
 C_z=TD_z .
\]
The map \(z\mapsto C_z\) is \cpc and order zero: it is the direct
sum of the order-zero maps \(\phi_i\), each multiplied by a
nonnegative scalar at most one.  It is supported on \(Q\), since
the taper \(T\) vanishes at level \(d\).
Put
\begin{equation}\label{norm:row-map}
 \lambda_0=(1+\eta^2)^{-1},\qquad
 V(z)=\lambda_0 BC_zB^*
     =\lambda_0\sum_{i,j}(1-l_j^{(i)}/d)
                    b_j^{(i)}\phi_i(z)(b_j^{(i)})^* .
\end{equation}
The last formula shows that \(V(z)\in A\), including when \(A\)
is nonunital.  Row compression gives complete positivity.
Furthermore \(0\le C_1\le1_{\mathcal L}\), so
\(\|V(1)\|\le\lambda_0\|B\|^2\le1\); hence \(V\) is contractive.

For orthogonal positive contractions \(z,w\in M_k\), we have
\(C_zC_w=0\).  Thus
\[
 V(z)V(w)
   =\lambda_0^2 BC_z(B^*B-1_{\mathcal L})C_wB^*,
\]
and \eqref{norm:row-gram} gives
\[
 \|V(z)V(w)\|
       \le\lambda_0^2(1+\eta^2)\eta^2
       \le\eta^2 .
\]
This proves \eqref{norm:order-zero-error}.

We next bound the commutator.  Let \(E_l\) be the projection onto
all slots of level \(l\), summed over the prototypes.  By
\eqref{norm:banded-model}, the selfadjoint contraction \(M_a\)
has only its level diagonal and its two adjacent-level diagonals.
The operators
\[
 M_{a,+}=\sum_{l=0}^{d-1}E_{l+1}M_aE_l,\qquad
 M_{a,-}=M_{a,+}^*
\]
each have norm at most one: their separate summands have
orthogonal initial spaces and orthogonal final spaces.  The
level diagonal commutes with \(T\), and
\[
 [M_a,T]=d^{-1}(M_{a,+}-M_{a,-}),\qquad
 \|[M_a,T]\|\le2/d .
\]
Within each prototype the entries of \(M_a\) are scalar and the
operator \(\phi_i(z)\) is the same in every slot.  Thus
\([M_a,D_z]=0\), and
\begin{equation}\label{norm:taper-bound}
 \|[M_a,C_z]\|\le2/d\qquad(\|z\|\le1).
\end{equation}
With \(R_a=aB-BM_a\), selfadjointness gives
\(R_a^*=B^*a-M_aB^*\), and a direct expansion yields
\[
 [a,V(z)]
 =\lambda_0\big(R_aC_zB^*+
                 B[M_a,C_z]B^*-BC_zR_a^*\big).
\]
Because \(C_z=QC_zQ\), \eqref{norm:row-residual} bounds both
outside error terms.  Consequently
\[
 \|[a,V(z)]\|
 \le\lambda_0\big(2\eta\|B\|+(2/d)\|B\|^2\big)
 \le2\eta+2/d .
\]
This proves \eqref{norm:central-error}; the estimate is uniform
over the unit ball of \(M_k\).

Finally fix \(\sigma\in\mathcal K(x,\gamma)\).  The retained cover
gives an index \(i\) with \(\|\sigma(s_i)\|=1\), so that
\(h_i=\sigma(\phi_i(1))\) has norm one by
\eqref{norm:cone-norm-one}.  Denote the represented support by
\(P_i^\sigma\), and put \(X=\sigma(x)\), \(H=\sigma(V(1))\).
The unique level-zero slot of the \(i\)-th model has representative
\(b_1^{(i)}=1\) and taper coefficient one.  Every summand of
\eqref{norm:row-map} at \(z=1\) is positive, so
\(H\ge\lambda_0 h_i\).  Also
\[
 P_i^\sigma X^2P_i^\sigma
 \ge(P_i^\sigma XP_i^\sigma)^2
 \ge(\gamma/4)^2P_i^\sigma,
\]
where the first inequality follows by expanding
\(P_i^\sigma X(1-P_i^\sigma)XP_i^\sigma\ge0\), and the second
uses \eqref{norm:positive-compression}.
Since \(h_i=P_i^\sigma h_iP_i^\sigma\),
\[
 \|Xh_iX\|
  =\|h_i^{1/2}X^2h_i^{1/2}\|
  \ge(\gamma/4)^2\|h_i\|
  =(\gamma/4)^2 .
\]
Positive domination and \(\|X\|\le1\) now imply
\[
 \|XH\|\ge\|XHX\|
       \ge\lambda_0\|Xh_iX\|
       \ge\gamma^2/32,
\]
because \(\eta<1\) gives \(\lambda_0\ge1/2\).
This proves \eqref{norm:finite-visibility}.
\end{proof}

\begin{proof}[Proof of Proposition~\ref{prop:norm}]
Choose increasing finite sets \(F_n\) of selfadjoint contractions
whose union is dense in the selfadjoint unit ball of \(A\).
Apply Lemma~\ref{norm:finite-stage} with
\(F=F_n\), \(d=n\), and \(\eta=(n+1)^{-1}\), and denote its maps
by \(\nu_n\).  Complete positivity and contractivity pass to the
ultrapower.  The error in \eqref{norm:order-zero-error} tends to
zero, so the resulting map is order zero.  The estimate in
\eqref{norm:central-error}, density, and
\(\|\nu_n(z)\|\le\|z\|\) imply that its range commutes with \(A\).
All maps have norm at most one, which supplies the bound required
to form their representing sequences.  The final estimate in
Lemma~\ref{norm:finite-stage} holds for every \(n\) and every
\(\sigma\in\mathcal K(x,\gamma)\), with the same constant
\(\gamma^2/32\).  This proves the stated simultaneous coordinate
quantifiers.
\end{proof}

\begin{remark}\label{norm:scope}
Proposition~\ref{prop:norm} controls irreducible representations of
the coordinate algebra \(A\).  It supplies the representation
estimates needed for later ideal comparisons; it makes no
assertion of fullness in \(\Aomega\cap A'\) or in its annihilator
quotient.
\end{remark}
